\section{Artificial Intelligence}

\par
Understanding how the term "Artificial Intelligence" (AI) was coined, one can
look at the way it is currently used nowadays, and connect how there isn't
really a consensus on "Intelligence" in AI implies. Even almost 70 years later
in 2025, people and researchers still debate what AI really means (as
highlighted in "On the Measure of Intelligence" by François
Chollet\cite{DBLP:journals/corr/abs-1911-01547}.

\subsection{Classification}

\par
When applied to different data sets, AI is often classified in the following
categories:

\begin{itemize}
	\item \textbf{Tabular data:} Machine Learning (ML)
	\item \textbf{Image data:} Artificial (Narrow) Intelligence (AI) or Computer
	      Cision
	\item \textbf{Text data:} Artificial (Narrow) Intelligence (AI) or Natural
	      Language Processing (NLP)
	\item \textbf{Image and text data:} Artificial General Intelligence (AGI)
\end{itemize}

\par
The reason for specifying "Narrow" in AI is to differentiate it from AGI and
theoretically ASI (Artificial Super Intelligence). Narrow AI referers to AI that
excels at specific tasks but cannot generalize or learn beyond its programming.
This is the reason why in the previous classification, "Narrow" was used to
specify for focusing on only one type of data. When an AI can understand and
learn across various domains, it can be classified as AGI.

\subsection{Issues with defining intelligence}

\par
In 1950, the Turing Test was proposed by Alan Turing to determine whether a
machine can exhibit intelligent behavior equivalent to, or indistinguishable
from a human:

\begin{quote}
	The new form of the problem can be described in terms of a game which
	we call the \textbf{imitation game}. It is played with three people, a man
	(A), a woman (B), and an interrogator (C) who may be of either sex. The
	interrigator stays in a room apart from the other two. The object of the game
	for the interrogator is to determine which of the other two is the man and
	which is the woman\cite{10.1093/mind/LIX.236.433}.
\end{quote}

\par
If one were to abide by the Turing Test as the definition of Intelligence, then
current day AI systems would be considered intelligent, as some of them can
successfully pass the Turing Test.

\par
The Loebner Prize, a competition that awards the most human-like AI through a
standard Turing Test, many flaws with the test have been highlighted. For
instance, many competitors focused on developing the AI for just the
competition, making it \textit{seem} intelligent, being able to fool the Turing
Test in the context of the competition, but not being truly intelligent.

\par
Another critique to the notion of machine intelligence is the "Lady Lovelace
objection":

\begin{quote}
	The analytical engine has no pretensions whatever to originate anything. It
	can do whatever we know how to order it to perform
\end{quote}

\par
It's understandable as to why pioneers in the field believed in how machines
would later evolve into being able to develop intelligence with sufficient
complexity.

\par
In 1820, Charles Babbage's "Difference Engine" was a complex enough
machine to perform basic arithmetic operations, nowadays with a much more
compact machine one can perform billions of the same operations performed by
the "Difference Engine" in a second. If one where to think of the brain as a
machine that performs billions of operations per second, then what follows is to
assume that a sufficiently complex machine should be able to operate in a
similar fashion, and therefore be intelligent. However, this assumption is based
on the idea that the brain \textit{is} a machine, which is still up for debate.

\par
A good analogy to understand the issue with defining intelligence in being a
"complex machine" is the Chinese Room argument by John Searle, in which the idea
is proposed where a complex enough machine could operate in a way in which it
"acts" and "seems" intelligent, but in reality it doesn't understand the subject
matter at all.

\begin{quote}
	\par

\end{quote}

\par

% TODO: Define equations for intelligence
Intelligence of system is over scope (sufficient case)

$$
	I^{{\theta}_T}_{IS,scope} = {Avg}_{T\in scope} \left[
	\omega_T\cdot\theta_T \sum_{C\in C_{{ur}_T^{\theta_T}}} \left[
	P_C \cdot \frac{GD_{IS,T,C}^{\theta_T}}{
	P_{IS,T}^{\theta_T}+E^{\theta_T}_{IS,T,C}}
	\right] \right]
$$

Intelligence of system is over scope (optimal case):

$$
	I^{opt}_{IS,scope} = {Avg}_{T\in scope} \left[
	\omega_{T,\Theta}\cdot\Theta \sum_{C\in C_{{ur}_T^{opt}}} \left[
	P_C \cdot \frac{GD_{IS,T,C}^{\Theta}}{
	P_{IS,T}^{\Theta}+E^{\Theta}_{IS,T,C}}
	\right] \right]
$$
